\SourceLeaf{22}
\[
[1]\qquad \xi=\int\!\int\frac{\mathrm{d}p.\mathrm{d}q.\cos p.\sin p}{Ar}
\]
\par
Likewise, by Theorem Six we obtain
\[
\begin{aligned}
[2]\qquad \xi&=-\int\!\int\frac{\mathrm{d}p.\mathrm{d}q.\sin p}{r^3}.(a-x)\bigl(\frac{(a-x)x}{AA}+\frac{(b-y)y}{BB}+\frac{(c-z)z}{CC}\bigr)
\end{aligned}
\]
\par
Finally, Theorem Four gives us
\[
\begin{aligned}
[3]\qquad &\int\!\int\frac{\mathrm{d}p.\mathrm{d}q.\sin p}{r^3}\bigl(\frac{(a-x)x}{AA}+\frac{(b-y)y}{BB}+\frac{(c-z)z}{CC}\bigr)=0\\
&\hspace{6em}\text{or }=-\frac{4\pi}{ABC}
\end{aligned}
\]
according as point $M$ lies outside or inside the body.

We now regard the quantities $A,B,C$ as particular values of three variables $\alpha,\beta,\gamma$ chosen as follows. The differences $\alpha\alpha-\beta\beta$ and $\alpha\alpha-\gamma\gamma$ are to be constant. Thus $\xi$ can be regarded as a function of the variables $\alpha,\beta,\gamma$, or rather of one of them. We denote simultaneous variations of the quantities $\xi,\alpha,\beta,\gamma$ by the sign $\delta$. It follows readily from equation $[1]$ that as $\alpha,\beta,\gamma$ increase without limit, $\xi$ decreases beyond all bounds, since clearly even the least value of $r$ increases beyond all bounds. We must therefore set $\xi=0$ for $\alpha=\infty$. Differentiating equation $[1]$ in the form
\[
\alpha\xi=\int\!\int\frac{\mathrm{d}p.\mathrm{d}q.\cos p.\sin p}{r}
\]
using $\delta$ as the sign of variation gives
\[
\alpha\delta\xi+\xi\delta\alpha=-\int\!\int\frac{\mathrm{d}p.\mathrm{d}q.\cos p.\sin p.\delta r}{rr}
\]
But we have
\[
\begin{aligned}
r\delta r&=-(a-x)\delta x-(b-y)\delta y-(c-z)\delta z\\
&=-(a-x)\cos p.\delta\alpha-(b-y)\sin p.\cos q.\delta\beta-(c-z)\sin p\sin q.\delta\gamma\\
&=-(a-x)x.\frac{\delta\alpha}{\alpha}-(b-y)y.\frac{\delta\beta}{\beta}-(c-z)z.\frac{\delta\gamma}{\gamma}\\
&=-\alpha\delta\alpha.\bigl(\frac{(a-x)x}{\alpha\alpha}+\frac{(b-y)y}{\beta\beta}+\frac{(c-z)z}{\gamma\gamma}\bigr)
\end{aligned}
\]
by virtue of $\alpha\delta\alpha-\beta\delta\beta=0$, $\alpha\delta\alpha-\gamma\delta\gamma=0$. Hence
\[
\alpha\delta\xi+\xi\delta\alpha=\delta\alpha.\int\!\int\frac{\mathrm{d}p.\mathrm{d}q.x\sin p}{r^3}\bigl(\frac{(a-x)x}{\alpha\alpha}+\frac{(b-y)y}{\beta\beta}+\frac{(c-z)z}{\gamma\gamma}\bigr)
\]
\par
Subtracting equation $[2]$ multiplied by $\delta\alpha$ from this, after replacing $A,B,C$ by $\alpha,\beta,\gamma$, gives
